\documentclass[11pt,letterpaper]{article}
\usepackage[lmargin=1in,rmargin=1in,bmargin=1in,tmargin=1in]{geometry}
\usepackage{../style}

\setlength{\parindent}{0ex}

% -------------------
% Content
% -------------------
\begin{document}

\pagenumbering{gobble}

% Recitation: Factoring, Roots, & Discriminants
\phantom{} \hfill {\bfseries \Large Recitation: Factoring, Roots, \& Discriminants Solutions} \hfill \phantom{} \par\vspace{0.3cm}

% Problem 1
\noindent {\bfseries Problem 1.} If any of the following polynomials factor, factor them completely:
	\begin{2enumerate}
	\item $x^2 - 10x - 24$
	\item $25x^3 - 4x$
	\item $x^2 - 7x - 6$
	\item $x^3 + 8$
	\item $4x^2 + 8x + 3$
	\item $x^2 + 1$
	\item $x^4 + x^3 - 2x^2$
	\item $3x^2 - x - 2$
	\item $x^3 - 2x^2 + 5x - 10$
	\item $32x^2 - 76x + 35$ [Hint: This has roots $\tfrac{5}{8}, \tfrac{7}{4}$.]
	\end{2enumerate} 

\noindent {\bfseries Solutions.}
	\begin{enumerate}[(a)]
	\item Because $-12 \cdot 2= -24$ and $-12 + 2= -10$, we have\dots
		\[
		x^2 - 10x - 24= (x - 12)(x + 2)
		\]
	
	\item Factoring out $x$, we have $25x^3 - 4x= x(25x^2 - 4)$. But the second term we recognize as a difference of perfect squares, $a^2 - b^2= (a - b)(a + b)$, with $a= 5x$ and $b= 2$. So, we have\dots
		\[
		25x^3 - 4x= x(25x^2 - 4)= x(5x - 2)(5x + 2)
		\]
	
	\item This does not factor (nicely). Observe the discriminant $D= b^2 - 4ac= (-7)^2 - 4(1)(-6)= 73$ is not a perfect square.
	
	\item This is a sum of perfect cubes: $a^3 + b^3= (a + b)(a^2 - ab + b^2)$, with $a= x$ and $b= 2$, so we have\dots
		\[
		x^3 + 8= (x + 2)(x^2 - 2x + 4)
		\]
	
	\item Observe $4 \cdot 3= 12$ and then observe $2 \cdot 6= 12$ and $2 + 6= 8$. So, we have\dots
		\[
		4x^2 + 8x + 3= (4x^2 + 2x) + (6x + 3)= 2x(2x + 1) + 3(2x + 1)= (2x + 1)(2x + 3)
		\]
	
	\item This is a sum of perfect squares---which \textit{never} factor over $\mathbb{R}$. The discriminant is also $D= b^2 - 4ac= 0^2 - 4(1)(1)= -4$, which indicates it does not factor. Over the complex numbers, it factors as $(x - i)(x + i)$.
	
	\item Factoring out the common $x^2$, we have $x^4 + x^3 - 2x^2= x^2(x^2 + x - 2)$. As $-1 \cdot 2= -2$ and $-1 + 2= 1$, we have $x^2 + x - 2= (x - 1)(x + 2)$. So, we have\dots
		\[
		x^4 + x^3 - 2x^2= x^2(x^2 + x - 2)= x^2(x - 1)(x + 2)
		\]
	
	\item Observe $3 \cdot -2= -6$, and observe $-3 \cdot 2= -6$ and $-3 + 2= -1$. So, we have\dots
		\[
		3x^2 - x - 2= (3x^2 - 3x) + (2x - 2)= 3x(x - 1) + 2(x - 1)= (x - 1)(3x + 2)
		\]
	
	\item Because of the four terms, we suspect this may factor by grouping (which we use when $a \neq 1$). Observe\dots
		\[
		x^3 - 2x^2 + 5x - 10= (x^3 - 2x^2) + (5x - 10)= x^2(x - 2) + 5(x - 2)= (x - 2)(x^2 + 5)
		\]
	
	\item Using the Fundamental Theorem of Algebra, we know a quadratic factors as $a(x - r_1)(x - r_2)$, where $r_1, r_2$ are the roots. So, we have\dots
		\[
		32x^2 - 76x + 35= 32 \left( x - \tfrac{5}{8} \right) \left( x - \tfrac{7}{4} \right)= 8 \left( x - \tfrac{5}{8} \right) \cdot 4 \left( x - \tfrac{7}{4} \right)= (8x - 5)(4x - 7)
		\]
	\end{enumerate}



\newpage



% Problem 2
\noindent {\bfseries Problem 2.} If $p(x)$ is a polynomial has roots $x= -1, -1, -1, 1, 3$ and $y$-intercept $5$, determine the polynomial. Also, determine the domain and range of $p(x)$. \par\vspace{0.5cm}

\noindent{\itshape By the Fundamental Theorem of Algebra, we know that a polynomial factors as $a(x - r_1)(x - r_2) \cdots (x - r_{n-1}) (x - r_n)$, where the $r_i$ are the roots and $a$ is the leading coefficient. So, using the given information, we have\dots
	\[
	p(x)= a \big(x - (-1) \big) \big(x - (-1) \big) \big(x - (-1) \big) (x - 1) (x - 3)= a (x + 1)^3 (x - 1)(x - 3)
	\]
But we know the $y$-intercept is $y= 5$, i.e. $y= 5$ when $x= 0$. But then\dots
	\[
	\begin{gathered}
	p(x)= a (x + 1)^3 (x - 1)(x - 3) \\
	p(0)= a (0 + 1)^3 (0 - 1)(0 - 3) \\
	5= 3a \\
	a= \tfrac{5}{3}
	\end{gathered}
	\]
Therefore, $p(x)= \tfrac{5}{3} (x + 1)^3 (x - 1)(x - 3)$.} \par\vspace{0.45cm}



% Problem 3
\noindent {\bfseries Problem 3.} Find all possible solutions to the following:
	\[
	\dfrac{5}{2 - x}= x
	\] 

	\[
	\begin{gathered}
	\dfrac{5}{2 - x}= x \\
	5= x(2 - x) \\
	5= 2x - x^2 \\
	x^2 - 2x + 5= 0 
	\end{gathered}
	\]

\noindent{\itshape Using the Quadratic Formula with $a= 1, b= -2, c=5$, we have\dots}
	
	\[
	x= \dfrac{-b \pm \sqrt{b^2 - 4ac}}{2a}= \dfrac{-(-2) \pm \sqrt{(-2)^2 - 4(1)5}}{2(1)}= \dfrac{2 \pm \sqrt{-16}}{2}= \dfrac{2 \pm i \sqrt{16}}{2}= \dfrac{2 \pm 4i}{2}= 1 \pm 2i
	\]



% Problem 4
\noindent {\bfseries Problem 4.} Explain why the polynomial $100x^2 - 400x + 100$ does not factor `nicely'. Then use the quadratic formula to factor the polynomial completely. \par\vspace{0.3cm}

\noindent{\itshape Observe $100x^2 - 400x + 100$ factors `nicely' if and only if $x^2 - 4x + 1$ factors `nicely' because $100x^2 - 400x + 100= 100(x^2 - 4x + 1)$. The polynomial $x^2 - 4x + 1$ has discriminant $D= b^2 - 4ac= (-4)^2 - 4(1)(1)= 12$, which is not a perfect square. So, $x^2 - 4x + 1$ does not factor `nicely.' Now, using the quadratic formula, we know that $x^2 - 4x + 1$ has roots\dots
	\[
	x= \dfrac{-b \pm \sqrt{b^2 - 4ac}}{2a}= \dfrac{-(-4) \pm \sqrt{(-4)^2 - 4(1)1}}{2(1)}= \dfrac{4 \pm \sqrt{12}}{2}= \dfrac{4 \pm 2 \sqrt{3}}{2}= 2 \pm \sqrt{3}
	\]
Then, using the Fundamental Theorem of Algebra, we have\dots
	\[
	100x^2 - 400x + 100= 100 \big(x - (2 + \sqrt{3}) \big) \big(x - (2 - \sqrt{3}) \big)
	\]
}

\end{document}